% Pursuing Stacks (cotes 134-2 to 134-8): archive cote 134-4, site page 90 (PDF page 91), typescript page 177.
% typescript. Draft compared: ps3.tex:4620-4694.
% AI-transcribed by Claude Opus 5.5 (claude-opus-5-5, adaptive thinking, effort high), Message Batches API, batch ps2 (msgbatch_01NVwr1Yuo8qJf5rykeGBXPK), 2026-11-12; not yet checked by eye.
\scanpage{134-4/90}
\mspage{177}

The obvious implications here are
\[ (i)\hand{$\Longrightarrow$}(i')\hand{$\Longrightarrow$}(ii) \quad , \]
for the last implication we only make use , besides A being $\underline{W}$-aspheric,
that the functors $(i_{/a})^*$ transform the projection $\hand{$\Delta$}_1 \to e$ in (Cat),
which is in $\underline{W}$ by L 3'), into a map in $\underline{W}_{(A_{/a})}$ , i.e. that the corres-%
ponding map in (Cat)
\[ A_{/axi^{\hand{*}}(\hand{$\Delta$}_1)} \longrightarrow A_{/a} \]
be in $\underline{W}$ , which by the final object axiom implying that $A_{/a}\to e$ is in
$\underline{W}$ , amounts to demanding that the first hand side is $\underline{W}$-aspheric i.e.
axi$^*(\hand{$\Delta$}_1)$ $\underline{W}_A$ -aspheric.

So we are left with prooving (ii)\hand{$\Rightarrow$}(i). By the criterion (iii)
of prop. 3 we know already (assuming (ii)) that A is a local $\underline{W}$-test ca-%
tegory, hence a $\underline{W}$-test category as A is $\underline{W}$-aspheric (cor.1); indeed
we can use I = i$^*(\hand{$\Delta$}_1)$ as a $\underline{W}_A$-aspheric interval, using the two canoni-%
cal sections deduced from the canonical sections of $\hand{$\Delta$}^1$ . The fact that
these are disjoint follows from the fact that i(a) non empty for any a in
A - we did not yet have to use the contractibility assumption on the
categories i(a). Thus we are reduced to proving that i$^*$ is model preser-%
ving - the same will then hold for the functors $i_{/a}$ (as required in
part b) in def.6), as the assumption (ii) is clearly stable under \xout{xxx}
restriction to the categories $A_{/a}$ . As we know already that $i_A$ is
model preserving (prop.2 (i)\hand{$\Rightarrow$}(ii)), all we have to do is to prove the
composition $f_i = i_A i^*$ is model preserving. But this was proved \inserted[typed]{yesterday} in the
key lemma of page 156. We're through \unsure{!}\xout{x}

\uline{Remark} The presentation will be maybe a little more elegant, if \unsure{We|we}
compplement the definition of a $\underline{W}$-test functor by the definition of a
\uline{local} \uline{$\underline{W}$-test} \uline{functor} , by which we mean that the induced functors
$i_{/a} : A_{/a} \longrightarrow M$ are weak test functors, period - which \xout{implies}\xout{xxxx} \inserted[typed]{means also}
\xout{\illegible{1}} that the following conditions hold :

a) A is a local $\underline{W}$-test category (def. 3) , i.e, the functors
$(i_{/a})^*$ : $M \to (A_{/a})^\wedge$ are all model preserving;
b) the functors $(i_{/a})^*$ (for a in A) are model preserving .

Thus it is clear that \xout{a} \inserted[typed]{if i is a} $\underline{W}$-test functor \inserted[typed]{it} is a local $\underline{W}$-test functor \inserted[typed]{such} \xout{xxx}
that moreover A is $\underline{W}$-aspheric. The converse is'nt clear in general, be-%
cause it is'nt clear that if A is a $\underline{W}$-test category and moreover all
functors $(i_{/a})^*$ are modelizing, than $i^*$ is equally modelizing. The
criterion (i') of theorem 1 shows however that this is so in the case
when (M,W) =((Cat),$\underline{W}$ ), and when we assume moreover the objects i(a)
contractible. \xout{On the other hand} We could now reformulate the theorem
\margin{90}
